% Options for packages loaded elsewhere
\PassOptionsToPackage{unicode}{hyperref}
\PassOptionsToPackage{hyphens}{url}
%
\documentclass[
  UTF8,
  zihao=5]{ctexart}
\usepackage{amsmath,amssymb}
\usepackage{lmodern}
\usepackage{iftex}
\ifPDFTeX
  \usepackage[T1]{fontenc}
  \usepackage[utf8]{inputenc}
  \usepackage{textcomp} % provide euro and other symbols
\else % if luatex or xetex
  \usepackage{unicode-math}
  \defaultfontfeatures{Scale=MatchLowercase}
  \defaultfontfeatures[\rmfamily]{Ligatures=TeX,Scale=1}
  \setmainfont[]{Times New Roman}
  \setsansfont[]{Arial}
  \setmonofont[]{Consolas}
\fi
% Use upquote if available, for straight quotes in verbatim environments
\IfFileExists{upquote.sty}{\usepackage{upquote}}{}
\IfFileExists{microtype.sty}{% use microtype if available
  \usepackage[]{microtype}
  \UseMicrotypeSet[protrusion]{basicmath} % disable protrusion for tt fonts
}{}
\makeatletter
\@ifundefined{KOMAClassName}{% if non-KOMA class
  \IfFileExists{parskip.sty}{%
    \usepackage{parskip}
  }{% else
    \setlength{\parindent}{0pt}
    \setlength{\parskip}{6pt plus 2pt minus 1pt}}
}{% if KOMA class
  \KOMAoptions{parskip=half}}
\makeatother
\usepackage{xcolor}
\IfFileExists{xurl.sty}{\usepackage{xurl}}{} % add URL line breaks if available
\IfFileExists{bookmark.sty}{\usepackage{bookmark}}{\usepackage{hyperref}}
\hypersetup{
  pdftitle={面向大语言模型推理的ACLGraph图执行技术综述},
  pdflang={zh-CN},
  hidelinks,
  pdfcreator={LaTeX via pandoc}}
\urlstyle{same} % disable monospaced font for URLs
\usepackage[margin=23mm,headheight=15pt]{geometry}
\usepackage{longtable,booktabs,array}
\usepackage{calc} % for calculating minipage widths
% Correct order of tables after \paragraph or \subparagraph
\usepackage{etoolbox}
\makeatletter
\patchcmd\longtable{\par}{\if@noskipsec\mbox{}\fi\par}{}{}
\makeatother
% Allow footnotes in longtable head/foot
\IfFileExists{footnotehyper.sty}{\usepackage{footnotehyper}}{\usepackage{footnote}}
\makesavenoteenv{longtable}
\usepackage{graphicx}
\makeatletter
\def\maxwidth{\ifdim\Gin@nat@width>\linewidth\linewidth\else\Gin@nat@width\fi}
\def\maxheight{\ifdim\Gin@nat@height>\textheight\textheight\else\Gin@nat@height\fi}
\makeatother
% Scale images if necessary, so that they will not overflow the page
% margins by default, and it is still possible to overwrite the defaults
% using explicit options in \includegraphics[width, height, ...]{}
\setkeys{Gin}{width=\maxwidth,height=\maxheight,keepaspectratio}
% Set default figure placement to htbp
\makeatletter
\def\fps@figure{htbp}
\makeatother
\setlength{\emergencystretch}{3em} % prevent overfull lines
\providecommand{\tightlist}{%
  \setlength{\itemsep}{0pt}\setlength{\parskip}{0pt}}
\setcounter{secnumdepth}{5}
\usepackage{float}
\usepackage{booktabs,longtable,array,caption,xurl,etoolbox}
\usepackage{setspace}
\usepackage{fancyhdr}
\setstretch{1.2}
\setlength{\parindent}{2em}
\setlength{\parskip}{0.18em}
\setlength{\emergencystretch}{3em}
\setlength{\LTpre}{0.7em}
\setlength{\LTpost}{0.7em}
\AtBeginEnvironment{longtable}{\fontsize{8.6}{11}\selectfont\setlength{\tabcolsep}{3pt}}
\AtBeginEnvironment{CSLReferences}{\fontsize{9}{11}\selectfont\setlength{\parskip}{0.3em}}
\captionsetup{font=small,labelfont=bf,labelsep=quad}
\ctexset{section={format=\large\bfseries,beforeskip=1.4em,afterskip=.6em},subsection={format=\normalsize\bfseries,beforeskip=1em,afterskip=.4em}}
\pagestyle{fancy}
\fancyhf{}
\fancyhead[C]{\small ACLGraph图执行技术综述}
\fancyfoot[C]{\thepage}
\renewcommand{\headrulewidth}{0.3pt}
\setlength{\headheight}{15pt}
\urlstyle{same}
\makeatletter
\renewcommand{\maketitle}{\begin{center}\LARGE\@title\end{center}\vspace{0.35em}}
\makeatother
\hypersetup{hidelinks,pdfauthor={},pdftitle={面向大语言模型推理的ACLGraph图执行技术综述}}
% Chinese language and typography are handled by ctex.
\ifLuaTeX
  \usepackage{selnolig}  % disable illegal ligatures
\fi
\newlength{\cslhangindent}
\setlength{\cslhangindent}{1.5em}
\newlength{\csllabelwidth}
\setlength{\csllabelwidth}{3em}
\newenvironment{CSLReferences}[2] % #1 hanging-ident, #2 entry spacing
 {% don't indent paragraphs
  \setlength{\parindent}{0pt}
  % turn on hanging indent if param 1 is 1
  \ifodd #1 \everypar{\setlength{\hangindent}{\cslhangindent}}\ignorespaces\fi
  % set entry spacing
  \ifnum #2 > 0
  \setlength{\parskip}{#2\baselineskip}
  \fi
 }%
 {}
\usepackage{calc}
\newcommand{\CSLBlock}[1]{#1\hfill\break}
\newcommand{\CSLLeftMargin}[1]{\parbox[t]{\csllabelwidth}{#1}}
\newcommand{\CSLRightInline}[1]{\parbox[t]{\linewidth - \csllabelwidth}{#1}\par}
\newcommand{\CSLIndent}[1]{\hspace{\cslhangindent}#1}

\title{面向大语言模型推理的ACLGraph图执行技术综述}
\usepackage{etoolbox}
\makeatletter
\providecommand{\subtitle}[1]{% add subtitle to \maketitle
  \apptocmd{\@title}{\par {\large #1 \par}}{}{}
}
\makeatother
\subtitle{ACLGraph Execution for Large Language Model Inference: A
Survey}
\author{}
\date{}

\begin{document}
\maketitle

\begin{figure}
\hypertarget{taxonomy-overview}{%
\centering
\includegraphics[width=1\textwidth,height=\textheight]{figures/taxonomy.pdf}
\caption{ACLGraph 图执行技术的五类代表工作。}\label{taxonomy-overview}
}
\end{figure}

\hypertarget{ux6458ux8981}{%
\subsection*{摘要}\label{ux6458ux8981}}
\addcontentsline{toc}{subsection}{摘要}

大语言模型解码反复执行相似任务，请求成员、序列长度和缓存映射却逐轮变化。本文按机制直接改变的执行决策，将图执行技术分为图构造与捕获、图选择与回退、回放状态管理、执行调度与资源分配、图生命周期管理五类。基于截至
2026 年 11 月 15 日纳入的 19
项研究及相关规范和工程实现，本文分别讨论昇腾直接工作与可借鉴的 GPU
工作。即使两轮选择同一张图，服务仍需刷新本轮参数、约束消费者的读取顺序，并维护底层资源的寿命。本文以地址、值、范围、依赖和请求归属检查这些条件，以关键路径和生命周期成本比较不同机制的代价。

\textbf{关键词：}
ACLGraph；大语言模型推理；运行时任务图；动态状态；调度；图资产恢复

\hypertarget{abstract}{%
\subsection*{Abstract}\label{abstract}}
\addcontentsline{toc}{subsection}{Abstract}

Language model serving repeatedly executes similar tasks while batch
membership, sequence lengths and cache mappings change. We classify
mechanisms by the execution decision they change: graph construction and
capture, graph selection and fallback, replay state management,
execution scheduling and resource allocation, and graph lifecycle
management. The review covers 19 studies available by November 15, 2026,
together with versioned specifications and implementations. Within each
class we distinguish direct Ascend evidence from GPU designs that
motivate comparisons. Address validity, value freshness, valid ranges,
dependency order and request ownership qualify reuse across classes.
Critical-path and lifecycle costs place steady-state execution and
cold-start results in their respective settings.

\textbf{Keywords:} ACLGraph; LLM inference; runtime task graph; dynamic
state; scheduling; graph restoration

\hypertarget{ux8303ux56f4ux4e0eux5206ux7c7bux65b9ux6cd5}{%
\subsection*{范围与分类方法}\label{ux8303ux56f4ux4e0eux5206ux7c7bux65b9ux6cd5}}
\addcontentsline{toc}{subsection}{范围与分类方法}

图回放复用已经记录的设备任务及其依赖。连续批处理允许请求在解码轮次之间加入和退出，分页
KV
缓存又让序列的逻辑位置与存储位置分离。即使批大小相同，下一轮要执行的任务、读取的数据和可用的资源也可能发生不同变化。{[}\protect\hyperlink{ref-yu2022orca}{1}--\protect\hyperlink{ref-zheng2024sglang}{3}{]}

图 1 以具体机制为分类单位。以 xLLM
为例，决定哪些模块能捕获属于第一类；为当前输入选择部分图或即时执行属于第二类；把变化的长度作为算子参数传入属于第三类。实验若同时启用这三项设计，其结果对应三者的组合。正文各章先给出归类判据，再分别讨论昇腾直接工作和
GPU
工作。前者包括昇腾系统设计、版本化接口、固定实现与部署观察，后者提供可比较的算法或系统选择。{[}\protect\hyperlink{ref-liu2025xllm}{4}{]}

\hypertarget{capture}{%
\section{图构造与捕获}\label{capture}}

\textbf{划分依据。}
改动决定一张图记录哪些设备任务、在哪个算子处分段，或用什么粒度表示任务与依赖，就归入本类。分类时可对照捕获前后的任务图：一次迁移即使保留了原算子，只要使原先逐轮单独提交的任务成为图节点，也改变了捕获边界。

\hypertarget{ux6607ux817eux76f4ux63a5ux5de5ux4f5c}{%
\subsection{昇腾直接工作}\label{ux6607ux817eux76f4ux63a5ux5de5ux4f5c}}

CANN
将捕获流上的任务记录在模型运行实例（modelRI）中，之后可以整体提交。多流捕获要求支流汇回主流。对于每轮需要单独执行的动态任务，运行时可以分别捕获其前后的静态段，在两个实例之间提交该任务。这样保留了每段内部的提交复用，代价是额外的实例、段间下发和同步。{[}\protect\hyperlink{ref-cann900overview}{5}--\protect\hyperlink{ref-cann910update}{8}{]}

xLLM 的 Adaptive Graph 展示了服务框架怎样选择捕获边界。FFN 和 LayerNorm
主要随有效 token
数变化，可以组成部分图；多头注意力还依赖批大小和查询、KV
长度，可保留即时执行路径。另一项昇腾实现把已有的 SAS/QLI
设备元数据任务纳入特定模型的完整解码图，使每轮主机准备和提交的边界随之变化。该改动适用于所列模型、算子与图模式。{[}\protect\hyperlink{ref-liu2025xllm}{4},\protect\hyperlink{ref-vllmascendpr21259}{9}{]}

\hypertarget{ux53efux501fux9274ux7684-gpu-ux5de5ux4f5c}{%
\subsection{可借鉴的 GPU
工作}\label{ux53efux501fux9274ux7684-gpu-ux5de5ux4f5c}}

IOS
在给定算子依赖图上搜索阶段划分，测量一个阶段合并或并行执行的代价。AutoGraph
先尝试等价图重写，再用实际运行筛选并行计划。两者都提示，决定捕获哪些任务之前，还可以改变候选任务结构。PyTorch
2 的 guards 和 graph breaks
则界定可编译的程序区域；把编译结果交给设备图回放，仍需检查设备捕获的约束。将这些思路用于昇腾时，需要按
CANN
接口和实际算子重新确认可捕获性。{[}\protect\hyperlink{ref-cann910single}{6},\protect\hyperlink{ref-ding2021ios}{10}--\protect\hyperlink{ref-ansel2024pytorch2}{12}{]}

Mirage 把优化对象推到代数与调度、内核的联合搜索：μGraph 统一
kernel/block/thread
三层，先以抽象表达剪枝、再做等价检查，最后做布局、顺序与内存规划
{[}\protect\hyperlink{ref-wu2025mirage}{13}{]}。原文 §7
说明当前实现只做一次不迭代的随机测试，并计划在流程末端为最佳 μGraph
增加提供理论保证的最终验证步骤
{[}\protect\hyperlink{ref-wu2025mirage}{13}{]}。理论保证、随机测试配置与浮点误差检查必须分别叙述；当前搜索过滤不能自动继承完整理论误差界。Mirage
也不提供 ACLGraph 运行时或跨进程图恢复，μGraph 不是 CUDA
Graph，片段收益不能直接替代完整模型收益。

\textbf{比较边界。}
可记录有效捕获区域、段数、图内任务和每轮图外提交次数。若同时做了融合或更换了内核，需保留这些变化各自的对照，才能解释捕获边界本身的作用。

\hypertarget{selection}{%
\section{图选择与回退}\label{selection}}

\textbf{划分依据。}
给定一组候选图，规则决定当前请求或批次命中哪张图、如何填充，或何时采用即时执行，就归入本类。切出一个可捕获的注意力外子图属于图构造；对某一批输入判定该子图是否可用，则属于图选择。

\hypertarget{ux6607ux817eux76f4ux63a5ux5de5ux4f5c-1}{%
\subsection{昇腾直接工作}\label{ux6607ux817eux76f4ux63a5ux5de5ux4f5c-1}}

vLLM-Ascend
的图包装器读取执行模式和批次描述符，决定复用已有图、捕获新图或即时执行。以
\href{https://github.com/vllm-project/vllm-ascend/pull/20279}{PR
\#20279} 为例，DFlash2 草稿在支持的完整图模式下使用
ACLGraph，注意力后端不兼容时回退到即时执行。{[}\protect\hyperlink{ref-ascend20279}{14}{]}SGLang
的 NPU
后端按形状持有图，并按算子约定填充和刷新槽位。容量准入只检查了一个条件：还要确认填充行不会改写真实请求的缓存。{[}\protect\hyperlink{ref-vllmascenddesign}{15}--\protect\hyperlink{ref-sglang2026impl}{18}{]}

xLLM
依据输入形状在部分图和即时执行路径间选择，并配合多图缓存和维度参数化。其
Table 8 报告了两个 Qwen3 模型在输入、输出长度均为 2048
时启用整套图模式的吞吐变化。Ascend 部署观察则记录了量化文本模型使用
full-decode 图、多模态负载使用 eager
的经历。前者的消融把几种机制一起打开；后者涉及不同负载。这两类证据都保留了各自的适用条件。{[}\protect\hyperlink{ref-liu2025xllm}{4},\protect\hyperlink{ref-yu2026ascendfield}{19}{]}

\hypertarget{ux53efux501fux9274ux7684-gpu-ux5de5ux4f5c-1}{%
\subsection{可借鉴的 GPU
工作}\label{ux53efux501fux9274ux7684-gpu-ux5de5ux4f5c-1}}

LightLLM 的 GPU
实现选择能容纳当前批次的最小可用桶，并检查最大上下文长度。FlashInfer
针对不同查询长度和格式配置准备多个图，运行时按当前 KV 配置分派。Hybrid
JIT
在出现新长度时捕获并缓存相应变体。三者分别展示了容量分桶、按配置选图和按需扩展候选图集合的取舍。{[}\protect\hyperlink{ref-lightllm2026code}{20}--\protect\hyperlink{ref-yadav2026hybridpreprint}{22}{]}

\textbf{比较边界。} 对同一请求序列同时记录实际命中模式、填充 token
数、回退频率和额外捕获成本。细分桶能减少某些批次的填充，也可能增加图实例和内存占用；共享模板或工作区会改变这一关系。{[}\protect\hyperlink{ref-vllmascendmode}{16},\protect\hyperlink{ref-ye2025flashinfer}{21},\protect\hyperlink{ref-liu2026foundry}{23}{]}

\hypertarget{replay-state}{%
\section{回放状态管理}\label{replay-state}}

\textbf{划分依据。}
选定图和实际任务结构之后，机制改变本轮参数或元数据由谁生成、存在哪里、怎样绑定、何时允许消费者读取，就归入本类。请求退出后即使复用同一块缓冲，旧请求的长度、索引和输出归属也须更新。比较每个方案时，应从生产者一直追到消费点。

\hypertarget{ux6607ux817eux76f4ux63a5ux5de5ux4f5c-2}{%
\subsection{昇腾直接工作}\label{ux6607ux817eux76f4ux63a5ux5de5ux4f5c-2}}

CANN
支持在更新流上更新匹配的任务组，组内任务数量与类型须符合规范。外部事件连接更新与图内消费，确保消费者读取对应状态版本。torch\_npu
在捕获时为已注册算子记录任务组和参数，并在回放时刷新。vLLM-Ascend 以
provider 解析本轮参数；SGLang 的 NPU 后端还用 fence
约束相邻轮次对同一存储的覆盖。各层分别承担接口、绑定和同步责任。{[}\protect\hyperlink{ref-cann910update}{8},\protect\hyperlink{ref-vllmascend2026impl}{17},\protect\hyperlink{ref-sglang2026impl}{18},\protect\hyperlink{ref-cann850updateapi}{24},\protect\hyperlink{ref-torchnpu2026graphs}{25}{]}

vLLM-Ascend
\href{https://github.com/vllm-project/vllm-ascend/pull/21191}{PR
\#21191} 处理的注意力元数据构建路径把 \texttt{query\_start\_loc\_cpu}
复制为新设备张量，而公共元数据已含设备边界；在合并的 MTP draft decode
图中，这一重复 H2D 会被捕获并每轮重放。修复方式是直接复用
\texttt{common\_attn\_metadata.query\_start\_loc{[}:\ num\_reqs\ +\ 1{]}}，使
FA3 收到的设备张量别名到 runner
维护的公共缓冲，而非独立副本；共享缓冲行为需跨消费者正确性验证
{[}\protect\hyperlink{ref-vllmascendpr21191}{26}{]}。该改动同时减少了
\(G\) 中的冗余拷贝任务，并将独立副本改为 \(M\)
中的共享别名；新的读写寿命责任随之出现。

另一项 draft 模型仅在 eager 时选择设备长度 CANN FlashAttn，启用 draft
ACL 图后回退到 FIA，而 FIA 的元数据构建器把精确设备序列长度物化为 Python
列表并在主机更新注意力任务。改动允许合规配置在 eager 与图模式下都用
FlashAttn，长度准备与 tiling 留在图捕获内，同一注意力组内兼容层通过既有
\texttt{get\_capture\_resource()} 按元数据身份、token
数与注意力几何共享生产者；图所有权提供缓存寿命，回放对持久、经拒绝校正的设备缓冲执行捕获的生产者，无需
FIA 任务更新
{[}\protect\hyperlink{ref-vllmascendpr21281}{27}{]}。这是``图资源共享''类机制：缓存键包含元数据身份与几何等兼容条件，图所有权限定复用窗口。

\hypertarget{ux53efux501fux9274ux7684-gpu-ux5de5ux4f5c-2}{%
\subsection{可借鉴的 GPU
工作}\label{ux53efux501fux9274ux7684-gpu-ux5de5ux4f5c-2}}

FlashInfer 把每轮规划和注意力执行拆成 plan 与 run。CPU
按序列长度计算工作队列及归约映射，写入 GPU 工作区的固定位置；可捕获的
run
从该位置读取当轮计划。多个层的计划规格匹配时可以共享结果。与昇腾图内元数据生产者对照，两者都让后续算子读取设备上的数据，生成数据的位置、传输开销和可共享条件却不同。{[}\protect\hyperlink{ref-ye2025flashinfer}{21},\protect\hyperlink{ref-vllmascendpr21281}{27}{]}

\begin{longtable}[]{@{}
  >{\raggedright\arraybackslash}p{(\columnwidth - 8\tabcolsep) * \real{0.20}}
  >{\raggedright\arraybackslash}p{(\columnwidth - 8\tabcolsep) * \real{0.20}}
  >{\raggedright\arraybackslash}p{(\columnwidth - 8\tabcolsep) * \real{0.20}}
  >{\raggedright\arraybackslash}p{(\columnwidth - 8\tabcolsep) * \real{0.20}}
  >{\raggedright\arraybackslash}p{(\columnwidth - 8\tabcolsep) * \real{0.20}}@{}}
\caption{动态请求的生成、存储与消费位置。生成位置与是否被捕获是独立维度；相同设备驻留状态可由不同位置的生产者形成。}\tabularnewline
\toprule
机制 & 变化如何表达 & 状态生成者／保存位置 & 消费与复用前提 &
证据类型 \\
\midrule
\endfirsthead
\toprule
机制 & 变化如何表达 & 状态生成者／保存位置 & 消费与复用前提 &
证据类型 \\
\midrule
\endhead
Orca/vLLM/SGLang原论文
{[}\protect\hyperlink{ref-yu2022orca}{1}--\protect\hyperlink{ref-zheng2024sglang}{3}{]}
& 批成员、块映射、多token工作量 & 服务调度与相应状态结构 &
不直接给ACLGraph协议 & 工作负载背景 \\
FlashInfer {[}\protect\hyperlink{ref-ye2025flashinfer}{21}{]} &
图族选择、CTA队列和归约映射 & CPU plan图外／固定device workspace &
run入图；匹配规格、容量及SM规划预算 & GPU系统与调度消融 \\
CANN任务更新 {[}\protect\hyperlink{ref-cann910update}{8}{]} &
同构任务组参数 & host更新流／运行实例任务 & 组内数量类型匹配；external
event & 版本化规范 \\
xLLM {[}\protect\hyperlink{ref-liu2025xllm}{4}{]} &
参数化、partial/eager、多图 & 配套算子与固定offset池 &
需算子改造；图模式组合消融 & 直接设计预印本 \\
MPK {[}\protect\hyperlink{ref-cheng2026mpk}{28}{]} &
图族选择、任务队列、MoE元数据 & 设备内运行时／设备内存 &
预生成tGraph与任务实现 & GPU预印本 \\
Ascend元数据生产者案例
{[}\protect\hyperlink{ref-vllmascendpr21259}{9},\protect\hyperlink{ref-vllmascendpr21281}{27}{]}
& 长度/tiling及SAS/QLI设备任务入图 & 设备计算／持久缓冲；捕获边界变化 &
模型、算子、图模式和共享键限制 & 局部实现证据 \\
\bottomrule
\end{longtable}

\textbf{比较边界。}
回放状态同时受地址、值、有效范围、依赖和请求归属约束。检查固定地址只能回答缓冲是否仍可访问；还可以通过请求退出、槽位复用和跨流覆盖等场景，在消费时分别核对值与归属。减少一项主机到设备的拷贝后，服务时延还取决于它是否原本位于关键路径。

\hypertarget{scheduling}{%
\section{执行调度与资源分配}\label{scheduling}}

\textbf{划分依据。}
在满足数据依赖的候选执行方式中，机制决定任务的提交顺序、流映射、并发窗口，或计算和通信能使用的资源，就归入本类。事件用于等待本轮元数据写好时，直接作用是保证状态可见；用同样的事件划出计算与通信可以重叠的区间时，研究对象转为执行调度。

\hypertarget{ux6607ux817eux76f4ux63a5ux5de5ux4f5c-3}{%
\subsection{昇腾直接工作}\label{ux6607ux817eux76f4ux63a5ux5de5ux4f5c-3}}

vLLM-Ascend
\href{https://github.com/vllm-project/vllm-ascend/pull/19716}{PR
\#19716} 在完整图路径中安排 WKV 投影与 MoE
专家分发重叠。另一项工程工作为注意力和打包操作按捕获容量安排并行资源。前者改变哪些任务可以同时运行，后者进一步决定并行时的资源分配。两项都是限定路径的实现证据；评价收益时，应同时观察各算子的持续时间和从分叉到汇合的总时间。{[}\protect\hyperlink{ref-ascend19716}{29},\protect\hyperlink{ref-ascend19909}{30}{]}

\hypertarget{ux53efux501fux9274ux7684-gpu-ux5de5ux4f5c-3}{%
\subsection{可借鉴的 GPU
工作}\label{ux53efux501fux9274ux7684-gpu-ux5de5ux4f5c-3}}

Nimble 为可并行的任务配置流和同步，Opara
在合法拓扑序内安排提交顺序，以减少资源碎片与并发干扰。NanoFlow
通过测量算子间的干扰来搜索流水和资源分配。POD-Attention
则在融合内核内按批组成调整 prefill 与 decode 的 CTA
分配。它们可分别用于比较合法并行、并发干扰和按负载调节资源三种问题。映射到昇腾时，需重新识别相应平台能控制的计算单元与通信资源。{[}\protect\hyperlink{ref-kwon2020nimble}{31}--\protect\hyperlink{ref-kamath2025pod}{34}{]}

Rammer、Mirage 与 MPK
代表另一条路线：不只调度既有算子，而是改变算子边界本身。Rammer 的
rOperator 暴露 rTask，vDevice/vEU 承载编译期 rProgram，Wavefront
策略联合选择 rKernel 与任务映射，CUDA 后端用 persistent thread blocks
实现
{[}\protect\hyperlink{ref-ma2020rammer}{35}{]}。其原文明确把机制与策略分离：Append/Wait
两个调度接口加剖析器构成机制，Wavefront
只是其中一个策略，作者预期该机制可支持更高级的策略
{[}\protect\hyperlink{ref-ma2020rammer}{35}{]}。因此 Rammer
的``最优''不是定理，而是特定 wave
启发式在剖析信息下的条件性选择；其静态条件是执行时间相对稳定的图与后端必须实现映射和同步。Opara
中的 Rammer− 仅是 Wavefront 策略的移植，不含完整 rTask/vDevice
系统，不能作为 Rammer 整体能力的证据
{[}\protect\hyperlink{ref-chen2025opara}{32}{]}。

MPK 与 Rammer 同样细化执行任务，并利用 Mirage
生成任务实现；其关注点进一步扩展到设备内执行模型：SM 级 tGraph
与事件、图规范化/线性化、mega-kernel 内的 worker 与 scheduler、混合
AOT/JIT 队列与分页共享内存
{[}\protect\hyperlink{ref-cheng2026mpk}{28}{]}。这里的JIT/AOT指任务入队时机，并非为任意新形状即时编译任务图。其原文对二者取舍给出明确机制解释：JIT
适应数据相关的执行时长不均（如 attention），但需 worker→scheduler→worker
两次同步；AOT
预入队只需一次事件同步，但无法吸收不均衡，故按算子是否数据相关分类，并在全局屏障后转回
AOT
{[}\protect\hyperlink{ref-cheng2026mpk}{28}{]}。这实际上是把``静态条件''从整图粒度细化到算子粒度：\(G\)
仍是按配置生成的静态 tGraph 族，\(\theta\) 中的批大小用于在多个 tGraph
间选择，而算子级的 JIT/AOT 标签是编译期生成的计划属性，复用它们须满足
\(\Phi\) 中的配置前提。MPK 的控制实验是离线批量推理（prompt 64、output
1024、最大批量 1--16），虽实现了 continuous batching 机制，但不是在线
SLO 验证；其 GPU SM/warp、TMA 与 NVSHMEM 假设不可直接套在 Ascend
上。截至本次核查，MPK 只按预印本引用，不能据致谢推断正式 venue。

\textbf{比较边界。}
并行策略可能让单个算子变慢，同时缩短整轮的最长依赖链。对照应固定输入与候选内核，记录跨流等待、资源竞争和整体时延。表
2 列出 GPU
相关路线在捕获或运行之前决定的对象，亦标明各自的静态前提和评测范围。

\begin{longtable}[]{@{}
  >{\raggedright\arraybackslash}p{(\columnwidth - 8\tabcolsep) * \real{0.20}}
  >{\raggedright\arraybackslash}p{(\columnwidth - 8\tabcolsep) * \real{0.20}}
  >{\raggedright\arraybackslash}p{(\columnwidth - 8\tabcolsep) * \real{0.20}}
  >{\raggedright\arraybackslash}p{(\columnwidth - 8\tabcolsep) * \real{0.20}}
  >{\raggedright\arraybackslash}p{(\columnwidth - 8\tabcolsep) * \real{0.20}}@{}}
\caption{计划生成与执行路线的横向比较。最优性和动态边界均相对于各自表示与搜索空间；本表不构成性能排名。}\tabularnewline
\toprule
路线 & 提前决定的对象 & 运行时仍变化的内容 & 主要成立条件 &
评测解释边界 \\
\midrule
\endfirsthead
\toprule
路线 & 提前决定的对象 & 运行时仍变化的内容 & 主要成立条件 &
评测解释边界 \\
\midrule
\endhead
Nimble {[}\protect\hyperlink{ref-kwon2020nimble}{31}{]} &
流映射、任务轨迹、预留内存 & 输入数据 & 静态网络、固定形状 &
V100推理/训练；逻辑并发目标 \\
IOS {[}\protect\hyperlink{ref-ding2021ios}{10}{]} & 阶段划分、合并或并发
& 固定配置下输入值 & 阶段代价可测量；有限图宽 & 四CNN；批量和设备专化 \\
AutoGraph {[}\protect\hyperlink{ref-zhao2023autograph}{11}{]} &
等价图替换与并行执行计划 & 固定图的输入值 & 给定规则、候选实测 &
七模型；启发式代价非全局证明 \\
Opara {[}\protect\hyperlink{ref-chen2025opara}{32}{]} &
流配置、资源感知提交序 & 固定图的输入值 & 剖析及资源类型有效 &
两GPU六模型；Rammer−非完整系统 \\
Rammer {[}\protect\hyperlink{ref-ma2020rammer}{35}{]} &
rKernel选择、rTask/vEU映射 & 计划内数据 & 可分解内核、后端同步 &
V100/MI50、初步IPU；覆盖不等同 \\
PyTorch 2 {[}\protect\hyperlink{ref-ansel2024pytorch2}{12}{]} &
guards保护的程序专化 & 受支持符号形状、分支回退 &
编译假设与图安全分别判断 & 多基准；不能把全部收益归图 \\
Mirage {[}\protect\hyperlink{ref-wu2025mirage}{13}{]} &
多层程序、布局和内存规划 & 专化输入数据 & 有限搜索与算子/验证范围 &
片段与集成；存在退化，基线已开图 \\
MPK {[}\protect\hyperlink{ref-cheng2026mpk}{28}{]} &
tGraph族、任务实现与标签 & 任务就绪、批量选择及路由元数据 &
CUDA设备任务和通信原语 & 预印本；离线批量而非在线SLO \\
\bottomrule
\end{longtable}

\hypertarget{lifecycle}{%
\section{图生命周期管理}\label{lifecycle}}

\textbf{划分依据。}
机制决定图实例何时准备、哪些工作区或代码需与实例一起保留、何时释放，以及环境变化后如何恢复，就归入本类。用缓存键决定本轮选哪张图属于第二类；用该键确定哪些实例常驻，以及失效时重建哪些资源，属于本类。

\hypertarget{ux6607ux817eux76f4ux63a5ux5de5ux4f5c-4}{%
\subsection{昇腾直接工作}\label{ux6607ux817eux76f4ux63a5ux5de5ux4f5c-4}}

CANN 的单流捕获规范要求相关资源在模型实例销毁前保持有效。xLLM 使用共享
HBM 池和固定偏移维持复用所需的地址关系。vLLM-Ascend
\href{https://github.com/vllm-project/vllm-ascend/pull/21373}{PR
\#21373}
在图预热期间准备静态内核；它改变的是准备工作的发生时机。评估这类改动应区分编译、捕获、首次可用时间及稳态回放。{[}\protect\hyperlink{ref-liu2025xllm}{4},\protect\hyperlink{ref-cann910single}{6},\protect\hyperlink{ref-ascend21373}{36}{]}

\hypertarget{ux53efux501fux9274ux7684-gpu-ux5de5ux4f5c-4}{%
\subsection{可借鉴的 GPU
工作}\label{ux53efux501fux9274ux7684-gpu-ux5de5ux4f5c-4}}

Medusa
在离线阶段记录任务、分配/释放轨迹和参数关系，以分配序列中的间接索引表示指针；在线阶段重放分配轨迹，将新返回的地址写回图节点。为处理同一地址被释放再分配造成的歧义，它把
kernel launch 与分配轨迹共同分析，而不是仅凭指针值匹配。其 kernel
恢复又结合名称查找、模块枚举与触发
kernel；重复层结构使第一层捕获能够触发后续所需模块。因此 Medusa
已恢复参数和 kernel
绑定，不能被描述为只保存拓扑，也不能把其方法整体简化为逐 kernel
手工补丁。{[}\protect\hyperlink{ref-zeng2025medusa}{37}{]}

Foundry 改变了恢复责任的放置：与其逐项把图中地址改成新地址，它用 VMM
重建图所预期的地址环境，并提取、重载所需二进制。两者都依赖可分析的执行环境，却选择了不同的稳定对象。Medusa
利用离线/在线控制流和分配次序的对应关系；Foundry
主动补记并重放捕获窗口分配，弥补跳过捕获造成的序列差异，同时固定KV容量。Foundry
的同拓扑模板还要求相关节点属性相容，跨rank复用需修补通信状态；PP不能直接套用所有rank共一模板的假设，需要按阶段处理。{[}\protect\hyperlink{ref-liu2026foundry}{23},\protect\hyperlink{ref-zeng2025medusa}{37}{]}

\textbf{比较边界。} Foundry
的计时以温进程、预热编译缓存与本地归档为前提，排除了权重加载和环境初始化。比较昇腾预热与
GPU
图恢复，应分别报告启动时间的起止点、所需的地址和代码兼容条件，以及复用或重建的具体资源。验证活动请求的
KV
数据、随机数状态与输出归属还需要另一套恢复协议。{[}\protect\hyperlink{ref-yu2026ascendfield}{19},\protect\hyperlink{ref-liu2026foundry}{23},\protect\hyperlink{ref-zeng2025medusa}{37}{]}

\hypertarget{ux7ed3ux8bed}{%
\section*{结语}\label{ux7ed3ux8bed}}
\addcontentsline{toc}{section}{结语}

五类机制可以组合。一个服务可以先分段捕获，按批次选图，在稳定地址上更新当轮参数，再按资源预算安排重叠，并在下次启动时重建兼容的实例。解释一次性能变化时，先定位哪项决策发生了变化，再检查正确性约束、代价所在的关键路径和实际测量的平台。下面的附录保留了资料范围、跨层分析框架、实验证据和仍待检验的研究问题。

\hypertarget{ux9644ux5f55-a-ux6587ux732eux8303ux56f4ux4e0eux8bc1ux636eux53e3ux5f84}{%
\section*{附录
A　文献范围与证据口径}\label{ux9644ux5f55-a-ux6587ux732eux8303ux56f4ux4e0eux8bc1ux636eux53e3ux5f84}}
\addcontentsline{toc}{section}{附录 A　文献范围与证据口径}

\hypertarget{ux95eeux9898ux5bfcux5411ux7684ux68c0ux7d22ux4e0eux7eb3ux6392}{%
\subsection*{问题导向的检索与纳排}\label{ux95eeux9898ux5bfcux5411ux7684ux68c0ux7d22ux4e0eux7eb3ux6392}}
\addcontentsline{toc}{subsection}{问题导向的检索与纳排}

本文资料截止于 2026 年 11 月 15
日。文献获取采用两组互补的定向检索：一组围绕运行时任务图、编译与 GPU
调度，另一组围绕 ACLGraph、NPUGraph、Ascend 推理及产生动态状态约束的 LLM
服务。检索使用论文题名、系统名、DOI 和组合关键词，例如``ACLGraph Ascend
inference paper NPUGraph''``FlashInfer dynamic planning CUDA
graph''及``Nimble Lightweight and Parallel GPU Task
Scheduling''；随后沿会议官网、作者全文、arXiv、厂商归档文档和固定源码提交核实方法与评测；恢复方向在交叉核查中补入
Medusa。该过程不是穷尽数据库的系统检索，不给出虚构的 PRISMA
流程或召回率。

纳入研究需要直接解释图计划、提交、调度、恢复，或明确产生服务图约束，并能取得相关方法和实验原文。官方文档和源码用于核实直接
ACLGraph 契约，PR 仅作局部机制案例。排除仅同名的知识图谱
nPUGraph、光感知 LightLLM 与 Arm Compute Library ACL
Graph；没有核实独立系统论文的 ModelTC LightLLM 只作为官方 GPU
实现参与比较。同一研究的预印本、发表稿和附录不重复计数。由此形成 19
项研究：2 项直接相关 Ascend 研究、14 项可迁移执行设计或编译研究、3
项工作负载背景研究；另以规范与实现补充软件栈映射。这里的数量描述本文证据集，不是对整个领域的规模判断。{[}\protect\hyperlink{ref-liu2025xllm}{4},\protect\hyperlink{ref-yu2026ascendfield}{19},\protect\hyperlink{ref-lightllm2026code}{20}{]}

证据池可分为四层。直接设计证据支持``提出了何种 ACLGraph
相关机制''，直接观察证据支持``特定部署遇到何种现象''；迁移设计支持跨平台机制比较，背景研究解释成员、寻址和程序状态为何变化。CANN
文档可支持接口语义，固定源码可支持责任和排序关系，却都不能独立支持服务收益。因此，后文表格同时保留机制、约束和验证范围，不把代码快照与同行评审论文当作同等强度的性能证据。

\hypertarget{ux7248ux672cux4e0eux7ecfux9a8cux7ed3ux8bbaux7684ux5bf9ux5e94}{%
\subsection*{版本与经验结论的对应}\label{ux7248ux672cux4e0eux7ecfux9a8cux7ed3ux8bbaux7684ux5bf9ux5e94}}
\addcontentsline{toc}{subsection}{版本与经验结论的对应}

版本身份决定可引用的结论。Opara 采用 2025
年正式卷期元数据，方法阅读来自作者接受稿；其 Rammer− 只是 Wavefront
策略移植。Mirage 采用 OSDI 2025 正式稿，MPK 则采用 2026 年更名后的 v2
预印本。Foundry 采用 2026 年 v1 预印本。Hybrid JIT 的 COMPSAC 2026
发表元数据已核实{[}\protect\hyperlink{ref-yadav2026hybrid}{38}{]}，但本文的方法与表格解读来自
arXiv
v1，未核实发表全文与其一致；因此后文明确限定数字的版本来源。{[}\protect\hyperlink{ref-wu2025mirage}{13},\protect\hyperlink{ref-yadav2026hybridpreprint}{22},\protect\hyperlink{ref-liu2026foundry}{23},\protect\hyperlink{ref-cheng2026mpk}{28},\protect\hyperlink{ref-chen2025opara}{32}{]}

直接实现按读取时的固定提交解释，不能以可移动的主分支名称代替版本。本文核对的
torch\_npu 提交为 2026 年 11 月 14 日快照，vLLM-Ascend 为 11 月 4
日快照，SGLang 为 11 月 15 日快照；后者在 11 月 2 日完成了独立 NPU
runner 向通用 runner/backend 的迁移。规范分别标明 CANN 9.0、9.1 或商业
8.5 文档路径，不把其能力与限制拼成一个无版本的``ACLGraph
最新规范''。具体读取范围和文件定位保留在文献登记中。{[}\protect\hyperlink{ref-vllmascend2026impl}{17},\protect\hyperlink{ref-sglang2026impl}{18},\protect\hyperlink{ref-cann850updateapi}{24},\protect\hyperlink{ref-torchnpu2026graphs}{25},\protect\hyperlink{ref-sglang2026refactor}{39}{]}

ACLGraph 的公开证据还具有特殊结构。xLLM 技术报告直接讨论 Ascend 上的
Adaptive Graph，并给出图模式开关消融；Ascend
实地研究记录了复杂负载部署中采用或放弃捕获的经历。前者是设计预印本，后者是观察性预印本。与此同时，CANN
规范、torch\_npu
和服务框架源码能够解释真实更新协议，但源码中的机制存在不等于跨配置性能成立。综述因此需要把学术设计、规范契约和实现案例连接起来，同时保持证据层次。{[}\protect\hyperlink{ref-liu2025xllm}{4},\protect\hyperlink{ref-vllmascend2026impl}{17}--\protect\hyperlink{ref-yu2026ascendfield}{19},\protect\hyperlink{ref-torchnpu2026graphs}{25}{]}

xLLM 的 Table 8 报告 Qwen3-1.7B 吞吐 2385.58→3038.88 tokens/s、Qwen3-4B
1539.93→1671.39，输入输出均为 2048
{[}\protect\hyperlink{ref-liu2025xllm}{4}{]}。但该表未单列卡数、请求率或重复次数，消融不能拆开各子机制贡献；xLLM
为预印本、非同行评审发表，其描述也未给出 CANN 逐 API
更新协议。因此这些数字只能作为``图模式开关在该配置下的观测''，不能作为各机制归因或跨配置外推的依据。

与之相对，Ascend field study 是观察性研究而非受控实验：16 卡 910、vendor
容器 0.13 系列，文本 W8A8 TP8×DP2 用 full-decode 图，VLM BF16 TP8×PP2 用
eager，二者并非 graph/eager 配对；作者因捕获问题使用 eager
并引用图资源警告
{[}\protect\hyperlink{ref-yu2026ascendfield}{19}{]}。该研究无图机制受控消融，不能把其他内核故障归因于图，旧容器观察也不代表
2026-11 最新能力。把 xLLM 的直接设计与 field study
的观察并置时，必须区分``系统提出了什么机制''与``现场观察到了什么现象''，xLLM
提供图模式开关的受控消融，但披露范围有限；实地研究提供非配对观察，二者具有不同的因果解释强度。

现有证据的可比性受四类差异制约：版本与发表状态、测量口径、硬件假设、以及基线实现程度。第一类差异在
Hybrid JIT 上最尖锐。其 COMPSAC 2026 的 DOI 与页码 196--204
已核实，但本次实读的是 arXiv
v1，未核对发表版的方法与表格是否一致；预印本文字称在全部评估生成长度上取得最低
P99，而其 Table II 显示 300--500 长度上 JIT+CUDA 为
17.31、18.90、20.50、22.09、23.68 ms，对应 TensorRT-LLM 为
16.54、16.53、16.54、16.61、16.65
ms，即长生成段并不占优{[}\protect\hyperlink{ref-yadav2026hybridpreprint}{22}{]}。这不是否定其短长度区间的收益，而是说明``全范围最低''的表述与自身表格矛盾，综述不能沿用该概括。其表I也存在算术范围问题：用各行基线除以Hybrid时延复算，HuggingFace范围约为1.85--2.90倍、TensorRT-LLM约为1.20--5.05倍，与正文给出的1.02--5.90及1.04--5.42不一致。复算只是文内一致性检查，不能替代重复实验；本文因此采用其机制和受限表格观察，不引用宽泛提升范围。{[}\protect\hyperlink{ref-yadav2026hybridpreprint}{22}{]}第二类差异在
xLLM：其图模式消融只覆盖两个模型、单一输入输出长度，且未披露卡数、请求率与重复次数，因此支持所报告条件下的图模式整体效果，但不能作为各子机制贡献的分解证据{[}\protect\hyperlink{ref-liu2025xllm}{4}{]}。第三类差异在
Ascend 现场研究：两个工作负载分别使用 full-decode 图与
enforce-eager，并非图与 eager
的配对对照，其并发甜点与冷启动观察属于部署观察，不能归因为图机制{[}\protect\hyperlink{ref-yu2026ascendfield}{19}{]}。第四类差异在
Foundry：计时排除权重加载与环境初始化，且假设兼容布局与二进制，其与
CUDA-checkpoint 的比较仅在单卡进行，因为后者不支持通信 kernel 所需的 IPC
内存{[}\protect\hyperlink{ref-liu2026foundry}{23}{]}。

MPK 的证据需要单独标注。其离线批量推理设置（prompt 64、output
1024、最大批量 1--16）支持 continuous batching 机制，但控制实验不是在线
SLO 验证；其 SM 级 tGraph、设备内 worker/scheduler、NVSHMEM
通信与寄存器按最大需求固定，都是 CUDA 特定假设，不能套用于
Ascend{[}\protect\hyperlink{ref-cheng2026mpk}{28}{]}。Mirage
的概率等价理论与其单次测试实现之间也存在距离：§7
说明当前为单次测试并计划添加最终保证步骤，因此不能把理论表述读作已执行的完整保证{[}\protect\hyperlink{ref-wu2025mirage}{13}{]}。这些限定不是要降低各工作的价值，而是要区分``机制已被证明存在''与``机制在目标条件下已被证明有效''。

\hypertarget{ux9644ux5f55-b-ux56feux8d44ux4ea7ux590dux7528ux6761ux4ef6ux4e0eux4ee3ux4ef7ux6a21ux578b}{%
\section*{附录
B　图资产、复用条件与代价模型}\label{ux9644ux5f55-b-ux56feux8d44ux4ea7ux590dux7528ux6761ux4ef6ux4e0eux4ee3ux4ef7ux6a21ux578b}}
\addcontentsline{toc}{section}{附录 B　图资产、复用条件与代价模型}

\hypertarget{ux56dbux79cdux5bf9ux8c61ux4e0eux4e0dux540cux7684ux590dux7528ux8fb9ux754c}{%
\subsection*{四种对象与不同的复用边界}\label{ux56dbux79cdux5bf9ux8c61ux4e0eux4e0dux540cux7684ux590dux7528ux8fb9ux754c}}
\addcontentsline{toc}{subsection}{四种对象与不同的复用边界}

模型语义图描述张量运算的数学含义；编译中间表示描述这些运算经过分解、融合、布局选择和专化后的程序；运行时任务图描述实际提交的
kernel、拷贝和同步操作；服务批次则是调度器在某一轮选择的请求与 token
集合。前三者是不同抽象层的程序表示，服务批次是执行工作集合，不必人为构造为另一张图。TorchDynamo
产生的 FX 图属于编译捕获接口，CUDA Graph 与 ACLGraph
属于运行时任务提交机制，二者同样使用``捕获''一词并不意味着捕获对象相同。{[}\protect\hyperlink{ref-cann900overview}{5},\protect\hyperlink{ref-ansel2024pytorch2}{12},\protect\hyperlink{ref-nvidia2025cudaguide}{40}{]}

各层之间也不是固定的一一映射。一个语义算子可能分解为多个
kernel，多个算子可能被融合为一个
kernel，一种服务批次又可能对应多个可用执行计划。Mirage 的 μGraph
将程序表示扩展到 kernel、block 和 thread 层次，MPK 的 tGraph
则把运行时依赖细化到 SM
级任务；它们改变的是表示与执行粒度，不能被等同为更大的 CUDA
Graph。服务层的请求成员变化也不必导致拓扑变化：如果算子能读取受约束的长度和索引，工作量可以在固定任务结构内变化。{[}\protect\hyperlink{ref-wu2025mirage}{13},\protect\hyperlink{ref-ye2025flashinfer}{21},\protect\hyperlink{ref-cheng2026mpk}{28}{]}

图 2
将这些对象与图资产连接起来。它强调两条关系：编译路径决定任务结构和代码，状态路径决定本轮选择、参数与数据。正确复用要求两条路径在消费点相容，不能仅检查编译产物存在，或仅检查服务批量落入某个桶。

\begin{figure}
\centering
\includegraphics[width=1\textwidth,height=\textheight]{figures/asset-contract.pdf}
\caption{图资产与动态状态的跨层连接。六元资产及五项不变量为本文综合框架；箭头表示职责关系，不代表厂商内部实现。}
\end{figure}

\hypertarget{ux516dux5143ux56feux8d44ux4ea7ux4e0eux53d8ux5316ux5206ux89e3}{%
\subsection*{六元图资产与变化分解}\label{ux516dux5143ux56feux8d44ux4ea7ux4e0eux53d8ux5316ux5206ux89e3}}
\addcontentsline{toc}{subsection}{六元图资产与变化分解}

本文把可复用执行所依赖的资产抽象为

\[ A=(G,\theta,M,B,C,\Phi). \]

\(G\) 为运行时任务拓扑及依赖；\(\theta\)
为节点参数模式与绑定，包括标量、指针参数和可更新入口；\(M\)
为地址映射、容量、布局及内存寿命；\(B\) 为 kernel
二进制和函数入口；\(C\)
为设备上下文、通信参与关系及其他执行环境；\(\Phi\)
为可复用前提集合。该表示统一描述本文比较问题，不声称六个分量能在任何实际
API 中独立更新。参数可能引用 \(M\) 中的地址，代码依赖
\(C\)，正是这种关联使拓扑相同仍不足以保证可执行。Foundry
对布局和二进制的恢复、TorchDynamo
对专化假设的检查，分别提供了这一综合的不同依据。{[}\protect\hyperlink{ref-ansel2024pytorch2}{12},\protect\hyperlink{ref-liu2026foundry}{23}{]}

第 \(t\) 轮的服务状态记为 \(s_t\)，包括请求成员、有效
token、长度、索引表、随机数状态和请求所有权。图选择函数将状态映射到候选资产
\(A_i\)，参数绑定过程再由 \(s_t\) 生成该次执行需要的 \(\theta_t\)
或设备可读数据。\(\Phi_i(s_t,C_t,M_t)\)
表示相关前提是否满足。它可以部分由静态类型或编译 guard
检查，部分由容量和配置准入保证，部分通过事件和资源所有权维护；不是要求每轮完整计算一个巨大的布尔谓词。这个划分允许比较检查发生的位置与成本，而不预设某一种契约语言。

变化可以进一步分解。输入值变化通常只改变缓冲内容；形状变化可能触发参数更新、图族选择或重新编译；控制流变化可能由图外分派、预建条件体或回退处理；任务时长变化即使不改变依赖，也可能影响静态资源分配的质量；进程、二进制或通信环境变化则影响资产能否恢复。请求退出和
slot 复用同时改变成员与所有权，不能简单归入 shape。FlashInfer、CUDA
条件节点、MPK 的混合队列和 Foundry
分别处理其中不同的变化，因而``是否支持动态''应由这些边界组成，而不能用单一布尔值表示。{[}\protect\hyperlink{ref-ye2025flashinfer}{21},\protect\hyperlink{ref-liu2026foundry}{23},\protect\hyperlink{ref-cheng2026mpk}{28},\protect\hyperlink{ref-nvidia2025cudaguide}{40}{]}

对于某个已选资产，静态性是相对于复用窗口而言的。例如，\(G\)
在多轮间不变，设备工作队列内容每轮可变；\(M\)
中虚拟地址稳定，映射到的逻辑请求却可以更换；\(B\)
在一个部署周期内不变，升级之后则需要重新判断兼容性。这里的``固定地址''首先指被消费者使用的有效地址及其约束，不应混同于物理页永不改变。vLLM
的分页机制已经将逻辑序列与物理存储分离，但其原论文并未自动建立 ACLGraph
的状态协议。{[}\protect\hyperlink{ref-kwon2023vllm}{2},\protect\hyperlink{ref-liu2026foundry}{23}{]}

\hypertarget{ux4e94ux9879ux4e0dux53d8ux91cfux53caux5176ux7ec4ux5408}{%
\subsection*{五项不变量及其组合}\label{ux4e94ux9879ux4e0dux53d8ux91cfux53caux5176ux7ec4ux5408}}
\addcontentsline{toc}{subsection}{五项不变量及其组合}

地址与寿命不变量要求：消费者使用的指针及其可访问范围，在执行期间引用有效资源；若图绑定了固定存储，则其生命周期必须覆盖相应执行。参数更新或间接寻址可以改变合法地址来源，因此该不变量不是``所有地址绝对不变''。CANN
单流捕获文档规定相关资源须保持至运行实例销毁，CUDA user objects
则提供异步资源引用管理，二者体现了不同接口下的寿命责任。{[}\protect\hyperlink{ref-cann910single}{6},\protect\hyperlink{ref-nvidia2025cudaguide}{40}{]}

值新鲜性不变量要求：消费者读到的是与本轮状态对应的值，而非捕获期或上一轮的值。同一地址可以存放正确的新长度，也可以残留旧长度，因而指针相等不能推出值正确。CANN
任务更新、torch\_npu
更新记录和服务框架参数绑定为部分参数提供了刷新路径；设备缓冲写入则是另一种路径。两者都需要继续核对更新对象与消费者是否匹配，不能用``存在
update
调用''替代这一关系。{[}\protect\hyperlink{ref-cann910update}{8},\protect\hyperlink{ref-vllmascend2026impl}{17},\protect\hyperlink{ref-sglang2026impl}{18},\protect\hyperlink{ref-torchnpu2026graphs}{25}{]}

范围不变量要求：有效请求数、有效 token 数、累计边界、上下文上界和
padding
语义相容，且不会越过所分配容量。它既约束资源上界，也约束有效前缀与填充尾部的含义。超过某个桶时，应选择能够覆盖输入的更大兼容桶、切换其他执行模式或重新构建；不能把输入交给更小容量图。FlashInfer
的上界布局和 LightLLM
的图准入是已有实例，但上界满足本身仍不能证明填充行不会误更新真实请求状态。{[}\protect\hyperlink{ref-lightllm2026code}{20},\protect\hyperlink{ref-ye2025flashinfer}{21}{]}

依赖不变量要求：生产者的状态发布先于相应消费者读取，当前读者完成又先于下一次可能冲突的覆盖。事件、流顺序和主机完成等待维护的是不同偏序边。额外全局屏障可能保守地保证正确性，却延长关键路径；没有必要的边则可能造成陈旧读取。CANN
普通事件用于跨流捕获传播，external event
用于图外更新与图内消费者之间的连接，二者不能无差别互换。{[}\protect\hyperlink{ref-cann910multi}{7},\protect\hyperlink{ref-cann910update}{8}{]}

请求世代不变量进一步要求：缓冲槽位、索引和输出属于当前请求及其有效执行轮次。这里的``世代''是本文对资源所有权更替的分析标签，不假定实现已有显式
generation 字段。请求甲退出后，请求乙可以合法复用同一
slot；地址、容量乃至长度都可能相同，但乙不能继承甲的过期索引和完成状态。vLLM
的块管理解释资源如何复用，SGLang 的 fence
解释相邻执行如何排序；将这些局部条件组合为请求语义仍是上层责任。{[}\protect\hyperlink{ref-kwon2023vllm}{2},\protect\hyperlink{ref-sglang2026impl}{18}{]}

五项不变量是相互关联的检查维度，而非本文已经证明的充要条件。范围错误可能造成无效地址，错误所有权可能表现为值错误；某些数值差异也可能来自浮点执行顺序，不能一律归因于世代或依赖问题。因此，对应的证据需要区分地址检查、消费时值检查、读写偏序和输出差分。局部断言能够排除某类错误，却不能自然升级为整个服务的正确性证明。

\hypertarget{ux4e09ux79cdux8fb9ux754cux60c5ux5f62}{%
\subsection*{三种边界情形}\label{ux4e09ux79cdux8fb9ux754cux60c5ux5f62}}
\addcontentsline{toc}{subsection}{三种边界情形}

请求退出首先改变有效工作集合。固定容量图可以保留原输入存储，把退出位置变成合法填充；也可以选择另一兼容图。两种方式都需要刷新长度、slot
与索引，而不是只减少主机列表长度。Orca
已解决迭代边界上不必等待整批完成的问题，但其选择性批处理不直接规定捕获图使用的稳定缓冲如何更新。由服务调度进入图执行，必须增加这一状态绑定环节。{[}\protect\hyperlink{ref-yu2022orca}{1},\protect\hyperlink{ref-ye2025flashinfer}{21}{]}

slot
复用把单轮合法性扩展为跨轮合法性。若甲的消费者仍在读共享元数据，乙提前覆写同一存储，即使每轮最终都能观察到正确值，也不能据此认定无竞争。图3只画出这一冲突别名集合上的写回、发布与消费顺序；独立host缓冲中的CPU规划、双缓冲或不相交设备槽中的生产可以提前进行，待合法时刻再绑定或拷入消费存储。FlashInfer区分pinned
host计划缓冲和device workspace，NanoFlow也将下一轮batch
formation与当前GPU执行重叠，说明状态生成不必等待所有当前读者。主机可先将replay入队，再完成更新；对应消费者须等待所读状态版本发布，不能仅凭Python行序判断可见性。{[}\protect\hyperlink{ref-cann910update}{8},\protect\hyperlink{ref-vllmascend2026impl}{17},\protect\hyperlink{ref-sglang2026impl}{18},\protect\hyperlink{ref-ye2025flashinfer}{21},\protect\hyperlink{ref-zhu2025nanoflow}{33}{]}

\begin{figure}
\centering
\includegraphics[width=1\textwidth,height=\textheight]{figures/state-publication.pdf}
\caption{同一共享存储上的冲突写回示例。读者完成约束同一别名集合的再覆写，发布连接对应状态版本的消费者；独立缓冲中的规划或生产可提前进行。}
\end{figure}

空 rank
则要求区分``没有本地有效工作''和``无需参与通信''。在需要共同参与的执行路径中，局部
token 数为零不自动授权跳过通信；元数据、填充及 collective
顺序必须与所用算法相容。这是分析性边界情形，不是对所有张量并行 rank
的工作分配描述。MPK 的计算通信任务表示、Foundry 的同结构跨 rank 模板以及
Ascend
元数据生产者案例分别提供相关机制；空工作可以在合法元数据协议内表示，也可能需要另外图族，不能预先断言其必然破坏模板或必然安全。{[}\protect\hyperlink{ref-vllmascendpr21259}{9},\protect\hyperlink{ref-liu2026foundry}{23},\protect\hyperlink{ref-cheng2026mpk}{28}{]}

\hypertarget{ux5173ux952eux8defux5f84ux4e0eux751fux547dux5468ux671fux4ee3ux4ef7}{%
\subsection*{关键路径与生命周期代价}\label{ux5173ux952eux8defux5f84ux4e0eux751fux547dux5468ux671fux4ee3ux4ef7}}
\addcontentsline{toc}{subsection}{关键路径与生命周期代价}

为避免把已重叠的操作重复计时，设第 \(t\)
轮实际执行的准备、更新、拷贝、计算、通信和等待构成依赖图
\(D_t\)。节点时长 \(d_v\)
应在给定资源分配及干扰条件下获得，并把导致等待的资源序列化关系纳入图中。以最长依赖链定义该轮的分析性时长：

\[ L_t=\max_{p\in\mathcal P(D_t)}\sum_{v\in p}d_v. \]

若尚未建模排队或资源干扰，\(L_t\)
只能是下界或结构解释，不是预测器。对同一执行，不应再把所有非关键路径节点的时长加回去；更新能否隐藏，取决于它是否推迟消费者以及与其他工作争用资源。

在部署窗口 \(\mathcal T\)
内，可以另行统计图构建、编译、更新、恢复和存储占用，并将目标写成条件化的选择问题：

\[ \min_{\pi,\mathcal A}\;\sum_{t\in\mathcal T}w_t L_t(\pi,\mathcal A)
 +\lambda_b K_{\mathrm{build}}+\lambda_r K_{\mathrm{restore}},
 \quad \text{s.t. } H(\mathcal A)\le H_{\max},\;R(\mathcal A)\le R_{\max},\;\Phi_i(s_t)=\mathrm{true}. \]

\(\mathcal A\) 是维护的图资产集合，\(\pi\) 是选择及更新策略，\(w_t\)
是工作负载权重，\(K\) 表示窗口内付出的构建或恢复工作，\(H\)
是保留内存，\(R\) 是实例、事件等资源用量，\(\lambda_b,\lambda_r\)
是分析所选权重。该式是本文提出的记账框架，未通过实验拟合；若构建或恢复已经进入
\(D_t\) 的前台关键路径，相同时间不能再次计入对应
\(K\)。它使稳态收益与资产维护成本可分开报告，并允许针对延迟、吞吐或启动期限采用不同目标。IOS
的搜索成本、MPK 的设备调度成本和 Foundry
的恢复成本由此能在共同生命周期中定位，而无需假定它们彼此可直接相加。{[}\protect\hyperlink{ref-ding2021ios}{10},\protect\hyperlink{ref-liu2026foundry}{23},\protect\hyperlink{ref-cheng2026mpk}{28}{]}

\hypertarget{ux8f6fux4ef6ux6808ux7684ux804cux8d23ux4e0eux66f4ux65b0ux987aux5e8f}{%
\subsection*{软件栈的职责与更新顺序}\label{ux8f6fux4ef6ux6808ux7684ux804cux8d23ux4e0eux66f4ux65b0ux987aux5e8f}}
\addcontentsline{toc}{subsection}{软件栈的职责与更新顺序}

更新有两条路径。方法一预先以动态任务为界分段捕获，前后任务存入不同model分别执行，动态任务在段间下发；并非每轮都需重捕获，适合大量任务更新但受硬件
model 数量限制；方法二以任务组 handle 标记、在更新流上以 begin/end
配对更新，适合少量单算子任务，且原文明确``更新任务比单独下发任务更耗时''，组内任务数量与类型必须匹配，任务组不支持多线程多流并发更新
{[}\protect\hyperlink{ref-cann910update}{8}{]}。并发更新场景需用
\texttt{ACL\_EVENT\_EXTERNAL} 事件连接更新流与图内消费者
{[}\protect\hyperlink{ref-cann910update}{8}{]}。8.5 的 API
页进一步限定更新流不可处于捕获状态，并按产品及 HDK
版本规定适用范围和可更新任务总量限制
{[}\protect\hyperlink{ref-cann850updateapi}{24}{]}。这些规范共同界定了
\(\Phi\)
的运行时下限：更新是显式的、有代价的、受资源与并发约束的，不能默认为廉价。

条件任务是另一条边界。\texttt{aclmdlRIAddCondTask} 在 9.1 中标注为试用
API、可能变更、不可用于生产环境，支持表为 A2/A3 推理产品及 950DT，Atlas
200I/500 A2 与 Atlas 推理/训练产品不支持
{[}\protect\hyperlink{ref-cann910conditional}{41}{]}。因此既不能说``CANN
普遍不支持条件''，也不能说``runner
已实现任意控制流''：底层条件任务入口存在，但其产品范围、试用状态与
runner 路径的采用是两件事，本次核查未在 runner
路径验证其采用、普适性或性能。

条件参数页进一步列出 IF、WHILE 与 SWITCH，以及分支数和 modelRI
数组表示；这证明规范已有受限条件与循环表示，而不是任意 Python
控制流。条件值的寿命同样重要：指定默认赋值标志时，每次执行前重置；不指定时由用户通过设备地址维护条件值。由此，控制选择本身也服从值新鲜性和跨轮依赖，而非加入条件节点后便自动正确。{[}\protect\hyperlink{ref-audit_cann_cond_params}{42},\protect\hyperlink{ref-audit_cann_cond_create}{43}{]}

torch\_npu 的 NPUGraph 把上述语义封装为 Python/C++ 两层。Python 侧
\texttt{\_GraphDispatchMode} 在捕获时对注册算子记录 dispatch
record：取当前流与 ExternalEvent、调用 handler 预处理、解析 schema、以
\texttt{graph\_task\_group\_begin/end} 记录任务组 handle、保存弱引用参数
{[}\protect\hyperlink{ref-torchnpu2026graphs}{25}{]}。更新时
\texttt{update\_capture\_record} 校验输入数量与记录数一致，在
update\_stream 上逐条 \texttt{graph\_task\_update\_begin}、改写
kwargs、重放算子、\texttt{graph\_task\_update\_end} 并记录事件
{[}\protect\hyperlink{ref-torchnpu2026graphs}{25}{]}。C++ 侧 NPUGraph 的
capture 路由内存池，replay 执行实例并处理 RNG
状态；\texttt{raw\_npu\_graph\_exec} 的文档明确直接执行会绕过 replay
的簿记（如 RNG
处理）{[}\protect\hyperlink{ref-torchnpu2026graphs}{25}{]}。这一层的关键限定是：注册表支持不等于所有算子自动更新；未注册算子走
passthrough，不产生 dispatch record；构建与驱动兼容性未在本次核查验证。

vLLM-Ascend 的设计文档把职责分为通用 dispatch 与 Ascend
wrapper/backend：descriptor 缓存、backend 参数 hook、ExternalEvent
与回放排序，并说明切图与尺寸覆盖受资源限制
{[}\protect\hyperlink{ref-vllmascenddesign}{15}{]}。模式指南进一步区分
ACLGraph、Npugraph\_ex 与 Xlite
是不同层次/路径，支持组合与默认项按版本变化
{[}\protect\hyperlink{ref-vllmascendmode}{16}{]}。源码层面，\texttt{ACLGraphWrapper}
按 descriptor 缓存图，FULL 常见路径在前序加屏障；不同分支的
update/replay 主机调用顺序可以相反
{[}\protect\hyperlink{ref-vllmascend2026impl}{17}{]}。\texttt{UpdatableGraph}
注册 handle/event/provider，\texttt{resolve\_tasks} 按 provider
解析当轮参数并断言数量匹配，\texttt{update} 在 update\_stream 上逐任务
apply，并对 paged attention 复用 workspace 缓存
{[}\protect\hyperlink{ref-vllmascend2026impl}{17}{]}。这构成一个明确的
provider 协议：图持有更新任务，provider 在每轮提供参数，数量断言是
\(\Phi\)
的一部分。该路径说明图所有权、参数绑定与事件发布如何组合，不能据数量断言单独推导参数含义正确。

SGLang 的 NPU 图支持经历了一个可观察的演进。issue 8030 建立
NPUGraphRunner 路线，并指出 FIA 的 host 参数需每轮
update，这是更新的直接动机
{[}\protect\hyperlink{ref-sglang8030}{44}{]}。issue 12194
以经验性条目列举地址、更新、事件、内存与资源约束，其中``sum/max
必同步''属过度概括，本文不采用该断言，也不外推其芯片资源数
{[}\protect\hyperlink{ref-sglang12194}{45}{]}。当前源码快照显示，NPU
full graph backend 按 shape 持图，静态 buffer 承载当轮状态；单 worker
异步 update、replay 与 fence 避免跨轮改写：先等前轮 fence，再 submit
更新、提交 replay、等 host update 结束并记录新 fence
{[}\protect\hyperlink{ref-sglang2026impl}{18}{]}。历史 NPUGraphRunner
类保留 host 长度与模型分流逻辑
{[}\protect\hyperlink{ref-sglang2026impl}{18}{]}。2026-11-02 的提交删除
\texttt{npu\_graph\_runner.py} 等文件并修改 backend 与通用 runner
{[}\protect\hyperlink{ref-sglang2026refactor}{39}{]}。这是 backend
迁移而非去掉能力：删除类名不等于删除 ACLGraph
能力，当前主分支快照也不是已发布版本，且本次未运行分布式测试；CPU
镜像不可用时仍有 D2H 路径。

\begin{longtable}[]{@{}
  >{\raggedright\arraybackslash}p{(\columnwidth - 8\tabcolsep) * \real{0.20}}
  >{\raggedright\arraybackslash}p{(\columnwidth - 8\tabcolsep) * \real{0.20}}
  >{\raggedright\arraybackslash}p{(\columnwidth - 8\tabcolsep) * \real{0.20}}
  >{\raggedright\arraybackslash}p{(\columnwidth - 8\tabcolsep) * \real{0.20}}
  >{\raggedright\arraybackslash}p{(\columnwidth - 8\tabcolsep) * \real{0.20}}@{}}
\caption{ACLGraph软件栈的责任与未自动保证的关系。源码快照不同，不代表同一联合运行环境。}\tabularnewline
\toprule
层次 & 主要责任 & 动态状态入口 & 同步/寿命机制 & 不能自动推出 \\
\midrule
\endfirsthead
\toprule
层次 & 主要责任 & 动态状态入口 & 同步/寿命机制 & 不能自动推出 \\
\midrule
\endhead
CANN
{[}\protect\hyperlink{ref-cann910single}{6}--\protect\hyperlink{ref-cann910update}{8}{]}
& 任务捕获、运行实例与更新 & 任务组、条件值 & 跨流汇合、external
event、资源寿命 & LLM slot语义或最优重叠 \\
torch\_npu {[}\protect\hyperlink{ref-torchnpu2026graphs}{25}{]} &
Python记录及C++执行环境 & 已注册handler与更新记录 & allocator
pool、事件、RNG簿记 & 所有算子都能自动更新 \\
vLLM-Ascend
{[}\protect\hyperlink{ref-vllmascenddesign}{15},\protect\hyperlink{ref-vllmascend2026impl}{17}{]}
& mode/descriptor适配、图持有任务 & backend hook与provider &
图所有权、事件及前轮屏障 & 数量断言等于请求语义证明 \\
SGLang
{[}\protect\hyperlink{ref-sglang2026impl}{18},\protect\hyperlink{ref-sglang2026refactor}{39}{]}
& 通用runner选图、NPU backend刷新 & 按算子协议的长度和静态buffer &
单worker update、rebind fence & runner迁移等于能力删除 \\
\bottomrule
\end{longtable}

表3总结四层之间的职责。CANN提供任务与同步原语，torch\_npu记录可更新算子并管理捕获环境，服务层把本轮请求状态绑定到图。vLLM-Ascend的provider登记与SGLang的自动记录、长度字段和fence属于不同的协议组织方式：前者显式关联图与参数提供者，后者在通用runner与NPU
backend之间划分设备策略。两者都需要确认值、有效范围、图所有权和跨轮覆盖顺序，接口形式本身不足以判定性能优劣。{[}\protect\hyperlink{ref-cann910update}{8},\protect\hyperlink{ref-vllmascend2026impl}{17},\protect\hyperlink{ref-sglang2026impl}{18},\protect\hyperlink{ref-torchnpu2026graphs}{25}{]}

源码快照尤其不能代替资源模型。主机提交replay后再update可以与图内事件配合；共享pool也必须服从实际并发的存活区间。是否减少端到端时延，应回到本附录的关键路径模型及实际资源竞争，而不能由类名、调用行数或一次成功捕获推导。编译侧的Npugraph\_ex、运行时ACLGraph和可选XliteGraph在文档中被分别列出，解释其组合时应继续保留层次。{[}\protect\hyperlink{ref-vllmascendmode}{16}--\protect\hyperlink{ref-sglang2026impl}{18}{]}

\hypertarget{ux9644ux5f55-c-ux5f85ux9a8cux8bc1ux7684ux7814ux7a76ux95eeux9898}{%
\section*{附录
C　待验证的研究问题}\label{ux9644ux5f55-c-ux5f85ux9a8cux8bc1ux7684ux7814ux7a76ux95eeux9898}}
\addcontentsline{toc}{section}{附录 C　待验证的研究问题}

以下问题沿用综述的比较基线。各项假设都须在相应硬件、固定版本与相同输入下接受检验。

\hypertarget{ux53efux7ec4ux5408ux7684ux72b6ux6001ux6709ux6548ux6027ux5951ux7ea6}{%
\subsection*{可组合的状态有效性契约}\label{ux53efux7ec4ux5408ux7684ux72b6ux6001ux6709ux6548ux6027ux5951ux7ea6}}
\addcontentsline{toc}{subsection}{可组合的状态有效性契约}

最接近这一问题的工作已经提供了实质机制。TorchDynamo 用 guards
检查程序专化假设，FlashInfer 以 plan/run
和容量上界连接动态计划与固定执行，torch\_npu 记录任务更新，两个 Ascend
服务框架再通过 provider、算子协议和 fence
绑定本轮状态。因此研究问题不是再增加一个更新
API，而是局部保证在共享元数据、slot 复用、图/eager
交替和空工作参与通信的条件下如何组合。某个图的参数数量正确，不说明每个参数属于正确请求；某个
buffer
地址正确，也不说明上一轮读者已经完成。{[}\protect\hyperlink{ref-ansel2024pytorch2}{12},\protect\hyperlink{ref-vllmascend2026impl}{17},\protect\hyperlink{ref-sglang2026impl}{18},\protect\hyperlink{ref-ye2025flashinfer}{21},\protect\hyperlink{ref-torchnpu2026graphs}{25}{]}

可检验的假设是：在算子显式声明输入范围、所有权与发布依赖的前提下，轻量世代标签和消费前检查能够检测仅靠地址及数量断言不能发现的陈旧状态，而无需每轮全设备同步。研究应先限定算子和图覆盖集合，构造请求替换、缩批后扩批和路径交替轨迹，比较现有断言、仅数值差分与契约检查的错误检出范围、误报和关键路径开销。故障注入可以改变更新对象、有效前缀或发布时序，但须保留真实消费者路径。若标签不能区分合法别名与错误继承，或必须依赖全同步才能成立，该假设即不成立。这里的研究贡献是可组合语义与代价，不是多运行若干测试。

还应区分数学等价、浮点可接受性、执行资源有效性和请求语义一致。Mirage
的理论与实现说明了前两者需要不同证据；图恢复后指针可用属于第三者；请求替换后输出归属属于第四者。任何跨层证明都必须给出覆盖的对象与信任边界，不能把有限随机测试或输出
token
一致性扩展为绝对保证。{[}\protect\hyperlink{ref-wu2025mirage}{13},\protect\hyperlink{ref-liu2026foundry}{23}{]}

\hypertarget{ux56feux65cfux7ef4ux62a4ux4e0eux8d44ux6e90ux8c03ux5ea6ux7684ux8054ux5408ux9009ux62e9}{%
\subsection*{图族维护与资源调度的联合选择}\label{ux56feux65cfux7ef4ux62a4ux4e0eux8d44ux6e90ux8c03ux5ea6ux7684ux8054ux5408ux9009ux62e9}}
\addcontentsline{toc}{subsection}{图族维护与资源调度的联合选择}

xLLM 已结合参数化、partial/eager和多图缓存，Hybrid
JIT使用按长度缓存与后台捕获，上游vLLM以模式及descriptor分派；FlashInfer进一步连接图族选择、每步CPU计划、固定workspace与SM数规划入口。IOS、Opara和NanoFlow则提供阶段测量、资源顺序和干扰感知优化。因此选图与资源适应本身均非空白。本方向关注资产集合的增删维护：请求分布漂移时，如何在实例/事件、保留内存和启动期限约束下决定构建、保留或恢复哪些资产，并计入模板及工作区共享改变的维护成本。{[}\protect\hyperlink{ref-liu2025xllm}{4},\protect\hyperlink{ref-ding2021ios}{10},\protect\hyperlink{ref-ye2025flashinfer}{21},\protect\hyperlink{ref-yadav2026hybridpreprint}{22},\protect\hyperlink{ref-chen2025opara}{32},\protect\hyperlink{ref-zhu2025nanoflow}{33},\protect\hyperlink{ref-vllmcudagraphdesign}{46}{]}

一个具体假设是：在这样的约束下，以预期关键路径收益减去维护成本来选择资产的策略，优于只按历史命中率保留图或只按
padding
比例增设桶的策略。其机制依据是，图命中频繁仍可能有高更新成本，罕见图也可能避免无法接受的尾部回退。可采用包含平稳、突变和周期分布的相同请求轨迹，比较固定桶、已有混合模式、带FlashInfer式选图/动态plan及资源入口的组合、命中率淘汰和联合选择；保持内核及总资源预算一致，分别统计执行、构建、回退和恢复的前台暴露时间。若收益只来自多占用内存或更换内核，便不能归于图族策略；若构建成本长期无法摊销，联合策略应退化为更保守的图集合。这使研究能被实际反驳，而不是承诺所有负载均受益。

\hypertarget{ux8ba1ux5212ux751fux6210ux4f4dux7f6eux4e0eux8de8ux6d88ux8d39ux8005ux5171ux4eab}{%
\subsection*{计划生成位置与跨消费者共享}\label{ux8ba1ux5212ux751fux6210ux4f4dux7f6eux4e0eux8de8ux6d88ux8d39ux8005ux5171ux4eab}}
\addcontentsline{toc}{subsection}{计划生成位置与跨消费者共享}

FlashInfer 已证明主机生成的计划可以跨匹配层复用，MPK
已在静态图族中执行设备内任务调度，Ascend
实现案例也已将部分元数据生产者纳入捕获。因而``首次把计划放到设备''或``首次合并逐层规划''都不是恰当的问题。剩余关系在于：逻辑长度相同而
attention
几何、tiling、工作区或更新语义不同时，哪些计划分量能共享，哪些必须由消费者重新解释。图所有权提供缓存寿命，却不能替代共享键的语义完整性。{[}\protect\hyperlink{ref-vllmascendpr21259}{9},\protect\hyperlink{ref-ye2025flashinfer}{21},\protect\hyperlink{ref-vllmascendpr21281}{27},\protect\hyperlink{ref-cheng2026mpk}{28}{]}

可检验假设为：将消费者不变的长度/分区信息与布局相关计划分离，能在限定算子集合内减少重复生产，并保持与独立规划的输出和依赖一致；但只有当减少的主机准备与传输超过新增设备工作和工作区竞争时，关键路径才会缩短。研究需要在同一后端、同一覆盖图和同一请求轨迹上，分别比较独立
host 更新、共享 host
计划、图内独立生产者与图内共享生产者。记录准备、拷贝、事件等待、设备生产和消费的依赖，而不是只读
API
次数。若不兼容布局迫使大部分计划重建，或设备生产挤占计算资源而延长消费路径，则应识别这一负收益区间。该方法同时解释何时保留
host 更新更合理。

\hypertarget{ux8de8ux56feux4e0eux591aux7b97ux5b50ux7684ux8d44ux6e90ux548cux72b6ux6001ux7ea6ux675f}{%
\subsection*{跨图与多算子的资源和状态约束}\label{ux8de8ux56feux4e0eux591aux7b97ux5b50ux7684ux8d44ux6e90ux548cux72b6ux6001ux7ea6ux675f}}
\addcontentsline{toc}{subsection}{跨图与多算子的资源和状态约束}

POD已经按当前批组成调整融合attention内的CTA分配策略，并运行时选择2或4
CTAs/SM；FlashInfer已经提供图族分派、动态plan及SM数入口，NanoFlow已经搜索多算子流水与干扰下的资源分配。MPK还显式承担设备调度器的资源成本。因此，按长度或prefill/decode比例选择局部重叠配置不构成这里的新问题；在Ascend复现这一选择只能首先视为迁移验证。{[}\protect\hyperlink{ref-ye2025flashinfer}{21},\protect\hyperlink{ref-cheng2026mpk}{28},\protect\hyperlink{ref-zhu2025nanoflow}{33},\protect\hyperlink{ref-kamath2025pod}{34}{]}

与上一方向的资产集合增删不同，这里固定已维护的图资产与内核集合，研究多个算子及图之间的组合可行性和切换成本。假设attention、GEMM与通信有可切换但受捕获接口限制的实现，且若干图共享有上限的workspace或状态版本缓冲：局部最快配置可能延长临时结果寿命，与另一图的写入冲突，或因kernel入口、launch属性不兼容而需要换图或合法更新。POD的kernel内自适应不负责决定这些跨资产绑定与存储寿命，单纯独立选图也不会自动给出它们的联合代价。NanoFlow式干扰搜索必须被纳入比较；不能把跨算子优化本身重新命名为空白。所读机制尚未共同给出上述捕获兼容、状态冲突和切换代价的统一处理，以下仅提出这一条件性组合问题。{[}\protect\hyperlink{ref-ye2025flashinfer}{21},\protect\hyperlink{ref-zhu2025nanoflow}{33},\protect\hyperlink{ref-kamath2025pod}{34}{]}

可检验假设是：在相同候选图资产、kernel实现集合和内存/实例预算下，将状态版本寿命、合法切换及多算子干扰同时纳入可行性与关键路径模型，能相对分阶段组合减少前台等待或反复重绑定。强基线应按算子及可用后端采用POD式局部自适应或FlashInfer式图选择/资源入口，再在相同可执行候选集上加入NanoFlow式资源搜索；这不假定两种attention后端已经合为同一kernel。各策略均可选择相同变体，也都必须做相同别名检查、必要同步与容量限制。比较``先局部优化再修复冲突''和联合选择，在小规模上用穷举可行组合校验搜索差距，并分别消融干扰模型及状态寿命约束。测量跨图等待、重绑定、切换与总关键路径，以未拟合工作量检验泛化。若组合基线已吸收这些条件而同样有效，或收益依赖基线不可用的额外内存、内核或未摊销的重捕获，假设即不成立。研究还须先确认平台暴露的真实控制粒度，不能假定SM数规划入口等价于Ascend资源隔离。

\hypertarget{ux6062ux590dux517cux5bb9ux6027ux4e0eux90e8ux5206ux5931ux6548ux5904ux7406}{%
\subsection*{恢复兼容性与部分失效处理}\label{ux6062ux590dux517cux5bb9ux6027ux4e0eux90e8ux5206ux5931ux6548ux5904ux7406}}
\addcontentsline{toc}{subsection}{恢复兼容性与部分失效处理}

Medusa已通过分配索引和kernel定位恢复绑定，Foundry已重建地址、代码及通信环境并复用模板；后者也比较了CUDA-checkpoint与CRIU的整个CUDA
context快照路线。这些对象须分开：图资产恢复不自动延续活请求的数据、RNG或输出归属。本方向因此限定在无在途请求且相关设备操作已经完成的服务静止点，考察图资产的部分失效；不提出活请求checkpoint或迁移方案。{[}\protect\hyperlink{ref-liu2026foundry}{23},\protect\hyperlink{ref-zeng2025medusa}{37}{]}

具体问题是在驱动、硬件和并行布局固定时，单独改变算子二进制或预留容量，哪些模板、参数模式及地址绑定仍可复用。假设为：在所需地址、模块和实例接口可见的范围内，显式记录这些分量的兼容依赖，可在拒绝不合法组合的同时减少相对全重建的工作。对照应包括保守全重捕获、原系统支持的同环境恢复，以及遇到不支持变化即全重建的现有恢复策略；不能强迫原系统执行非法资产来制造优势。逐项改变条件并独立校验判定、初始化后的输出与依赖，记录图构建、加载、参数准备及前台时间，不用这些指标代表活请求恢复。仅拓扑检查可作为错误判定对照，不应让已知不兼容图实际执行。若变化迫使代码、容量和绑定整体失效，部分重建无收益亦是有效结果。ACLGraph是否具备所需可见原语须单独核实；研究贡献应落在兼容判据和失效传播，而不是补齐平台测量。

\hypertarget{ux53c2ux8003ux6587ux732e}{%
\section*{参考文献}\label{ux53c2ux8003ux6587ux732e}}
\addcontentsline{toc}{section}{参考文献}

\hypertarget{refs}{}
\begin{CSLReferences}{0}{0}
\leavevmode\vadjust pre{\hypertarget{ref-yu2022orca}{}}%
\CSLLeftMargin{{[}1{]} }%
\CSLRightInline{G.-I. Yu, J. S. Jeong, G.-W. Kim, 等. {Orca: A
Distributed Serving System for Transformer-Based Generative Models}.
16th USENIX Symposium on Operating Systems Design and Implementation,
2022, 521--538.
\url{https://www.usenix.org/conference/osdi22/presentation/yu}.}

\leavevmode\vadjust pre{\hypertarget{ref-kwon2023vllm}{}}%
\CSLLeftMargin{{[}2{]} }%
\CSLRightInline{W. Kwon, Z. Li, S. Zhuang, 等. {Efficient Memory
Management for Large Language Model Serving with PagedAttention}.
Proceedings of the 29th Symposium on Operating Systems Principles, 2023,
611--626. doi:
\href{https://doi.org/10.1145/3600006.3613165}{10.1145/3600006.3613165}.}

\leavevmode\vadjust pre{\hypertarget{ref-zheng2024sglang}{}}%
\CSLLeftMargin{{[}3{]} }%
\CSLRightInline{L. Zheng, L. Yin, Z. Xie, 等. {SGLang: Efficient
Execution of Structured Language Model Programs}. Advances in Neural
Information Processing Systems, 2024, 37, 62557--62583. doi:
\href{https://doi.org/10.52202/079017-2000}{10.52202/079017-2000}.}

\leavevmode\vadjust pre{\hypertarget{ref-liu2025xllm}{}}%
\CSLLeftMargin{{[}4{]} }%
\CSLRightInline{T. Liu, T. Peng, P. Yang, 等. {xLLM Technical Report}.
arXiv preprint, 2510.14686v2 (2026; first posted 2025), 2025. doi:
\href{https://doi.org/10.48550/arXiv.2510.14686}{10.48550/arXiv.2510.14686}.}

\leavevmode\vadjust pre{\hypertarget{ref-cann900overview}{}}%
\CSLLeftMargin{{[}5{]} }%
\CSLRightInline{Huawei. {ACL Graph Overview}.
\url{https://www.hiascend.com/doc_center/source/en/CANNCommunityEdition/900/programug/acldevg/runtime_doc_dev_0045.html}.}

\leavevmode\vadjust pre{\hypertarget{ref-cann910single}{}}%
\CSLLeftMargin{{[}6{]} }%
\CSLRightInline{Huawei. {Single-Stream Capture}.
\url{https://www.hiascend.com/doc_center/source/en/CANNCommunityEdition/910/others/acldevg/runtime_doc_dev_0030.html}.}

\leavevmode\vadjust pre{\hypertarget{ref-cann910multi}{}}%
\CSLLeftMargin{{[}7{]} }%
\CSLRightInline{Huawei. {Cross-Stream Capture}.
\url{https://www.hiascend.com/doc_center/source/en/CANNCommunityEdition/910/others/acldevg/runtime_doc_dev_0031.html}.}

\leavevmode\vadjust pre{\hypertarget{ref-cann910update}{}}%
\CSLLeftMargin{{[}8{]} }%
\CSLRightInline{Huawei. {Task Update}.
\url{https://www.hiascend.com/doc_center/source/en/CANNCommunityEdition/910/others/acldevg/runtime_doc_dev_0032.html}.}

\leavevmode\vadjust pre{\hypertarget{ref-vllmascendpr21259}{}}%
\CSLLeftMargin{{[}9{]} }%
\CSLRightInline{vLLM Ascend contributors. {{[}Performance{]}{[}Model{]}
Capture DeepSeek V4 decode metadata in full graphs}. 2026.
\url{https://github.com/vllm-project/vllm-ascend/pull/21259}.}

\leavevmode\vadjust pre{\hypertarget{ref-ding2021ios}{}}%
\CSLLeftMargin{{[}10{]} }%
\CSLRightInline{Y. Ding, L. Zhu, Z. Jia, 等. IOS: Inter-Operator
Scheduler for CNN Acceleration. Proceedings of Machine Learning and
Systems 3, 2021, 3, 167--180.
\url{https://proceedings.mlsys.org/paper_files/paper/2021/hash/1f8053a67ec8e0b57455713cefdd8218-Abstract.html}.}

\leavevmode\vadjust pre{\hypertarget{ref-zhao2023autograph}{}}%
\CSLLeftMargin{{[}11{]} }%
\CSLRightInline{Y. Zhao, Q. Sun, Z. He, 等. AutoGraph: Optimizing DNN
Computation Graph for Parallel GPU Kernel Execution. Proceedings of the
AAAI Conference on Artificial Intelligence, 2023, 37, 11354--11362. doi:
\href{https://doi.org/10.1609/aaai.v37i9.26343}{10.1609/aaai.v37i9.26343}.}

\leavevmode\vadjust pre{\hypertarget{ref-ansel2024pytorch2}{}}%
\CSLLeftMargin{{[}12{]} }%
\CSLRightInline{J. Ansel, E. Yang, H. He, 等. PyTorch 2: Faster Machine
Learning Through Dynamic Python Bytecode Transformation and Graph
Compilation. Proceedings of the 29th ACM International Conference on
Architectural Support for Programming Languages and Operating Systems,
Volume 2, 2024, 929--947. doi:
\href{https://doi.org/10.1145/3620665.3640366}{10.1145/3620665.3640366}.}

\leavevmode\vadjust pre{\hypertarget{ref-wu2025mirage}{}}%
\CSLLeftMargin{{[}13{]} }%
\CSLRightInline{M. Wu, X. Cheng, S. Liu, 等. Mirage: A Multi-Level
Superoptimizer for Tensor Programs. 19th USENIX Symposium on Operating
Systems Design and Implementation (OSDI 25), 2025, 21--38.
\url{https://www.usenix.org/conference/osdi25/presentation/wu-mengdi}.}

\leavevmode\vadjust pre{\hypertarget{ref-ascend20279}{}}%
\CSLLeftMargin{{[}14{]} }%
\CSLRightInline{vLLM Ascend contributors. {Add DFlash2 ACL graph}.
\url{https://github.com/vllm-project/vllm-ascend/pull/20279}.}

\leavevmode\vadjust pre{\hypertarget{ref-vllmascenddesign}{}}%
\CSLLeftMargin{{[}15{]} }%
\CSLRightInline{vLLM Ascend contributors. {ACL Graph design}.
\url{https://docs.vllm.ai/projects/ascend/en/v0.23.0/developer_guide/Design_Documents/ACL_Graph.html}.}

\leavevmode\vadjust pre{\hypertarget{ref-vllmascendmode}{}}%
\CSLLeftMargin{{[}16{]} }%
\CSLRightInline{vLLM Ascend contributors. {Graph Mode Guide}.
\url{https://docs.vllm.ai/projects/ascend/en/v0.23.0/user_guide/feature_guide/graph_mode.html}.}

\leavevmode\vadjust pre{\hypertarget{ref-vllmascend2026impl}{}}%
\CSLLeftMargin{{[}17{]} }%
\CSLRightInline{vLLM Ascend contributors. {ACLGraphWrapper and
UpdatableGraph}.
\url{https://github.com/vllm-project/vllm-ascend/tree/fef7a9886cf1365ae951ff01fb1599802ca5cc6a}.}

\leavevmode\vadjust pre{\hypertarget{ref-sglang2026impl}{}}%
\CSLLeftMargin{{[}18{]} }%
\CSLRightInline{SGLang contributors. {SGLang NPU full graph backend and
generic decode runner}.
\url{https://github.com/sgl-project/sglang/tree/f4c16239a557734690ff6652e392f36ac3788798}.}

\leavevmode\vadjust pre{\hypertarget{ref-yu2026ascendfield}{}}%
\CSLLeftMargin{{[}19{]} }%
\CSLRightInline{Z. Yu. {On the Limitations of Non-GPU AI Accelerators
for Large-Model Inference: A Field Study of MoE and Multimodal Serving
on Huawei Ascend}. arXiv preprint, 2607.08215v1, 2026. doi:
\href{https://doi.org/10.48550/arXiv.2607.08215}{10.48550/arXiv.2607.08215}.}

\leavevmode\vadjust pre{\hypertarget{ref-lightllm2026code}{}}%
\CSLLeftMargin{{[}20{]} }%
\CSLRightInline{ModelTC LightLLM contributors. {LightLLM decode/prefill
CUDA graph implementation}.
\url{https://github.com/ModelTC/LightLLM/tree/1a2a306d3dd8b245ced0830100a6d0b0d302407a}.}

\leavevmode\vadjust pre{\hypertarget{ref-ye2025flashinfer}{}}%
\CSLLeftMargin{{[}21{]} }%
\CSLRightInline{Z. Ye, L. Chen, R. Lai, 等. {FlashInfer: Efficient and
Customizable Attention Engine for LLM Inference Serving}. Proceedings of
Machine Learning and Systems, 2025, 7.
\url{https://proceedings.mlsys.org/paper_files/paper/2025/hash/dbf02b21d77409a2db30e56866a8ab3a-Abstract-Conference.html}.}

\leavevmode\vadjust pre{\hypertarget{ref-yadav2026hybridpreprint}{}}%
\CSLLeftMargin{{[}22{]} }%
\CSLRightInline{D. K. Yadav 和 T. Zhao. {Hybrid JIT-CUDA Graph
Optimization for Low-Latency Large Language Model Inference}. arXiv
preprint, 2604.23467v1, 2026. doi:
\href{https://doi.org/10.48550/arXiv.2604.23467}{10.48550/arXiv.2604.23467}.}

\leavevmode\vadjust pre{\hypertarget{ref-liu2026foundry}{}}%
\CSLLeftMargin{{[}23{]} }%
\CSLRightInline{X. Liu, Y. Wu, Y. Yao, 等. Foundry: Template-Based CUDA
Graph Context Materialization for Fast LLM Serving Cold Start. arXiv
preprint, 2604.06664v1, 2026. doi:
\href{https://doi.org/10.48550/arXiv.2604.06664}{10.48550/arXiv.2604.06664}.}

\leavevmode\vadjust pre{\hypertarget{ref-cann850updateapi}{}}%
\CSLLeftMargin{{[}24{]} }%
\CSLRightInline{Huawei. {aclmdlRICaptureTaskUpdateBegin}.
\url{https://www.hiascend.com/doc_center/source/en/canncommercial/850/API/appdevgapi/aclcppdevg_03_1825.html}.}

\leavevmode\vadjust pre{\hypertarget{ref-torchnpu2026graphs}{}}%
\CSLLeftMargin{{[}25{]} }%
\CSLRightInline{Ascend PyTorch contributors. {torch\_npu NPUGraph Python
and C++ implementation}.
\url{https://github.com/Ascend/pytorch/tree/6133dfa4ea363aa58e04609f4d5e3e366a4df17c}.}

\leavevmode\vadjust pre{\hypertarget{ref-vllmascendpr21191}{}}%
\CSLLeftMargin{{[}26{]} }%
\CSLRightInline{vLLM Ascend contributors.
{{[}Performance{]}{[}Attention{]} Reuse device query boundaries to avoid
redundant H2D}. 2026.
\url{https://github.com/vllm-project/vllm-ascend/pull/21191}.}

\leavevmode\vadjust pre{\hypertarget{ref-vllmascendpr21281}{}}%
\CSLLeftMargin{{[}27{]} }%
\CSLRightInline{vLLM Ascend contributors.
{{[}Performance{]}{[}SpecDecode{]} Enable FlashAttn for GLM DFlash2 ACL
graphs}. 2026.
\url{https://github.com/vllm-project/vllm-ascend/pull/21281}.}

\leavevmode\vadjust pre{\hypertarget{ref-cheng2026mpk}{}}%
\CSLLeftMargin{{[}28{]} }%
\CSLRightInline{X. Cheng, Z. Zhang, Y. Zhou, 等. MPK: A Compiler and
Runtime for Mega-Kernelizing Tensor Programs. arXiv preprint,
2512.22219v2, 2026. doi:
\href{https://doi.org/10.48550/arXiv.2512.22219}{10.48550/arXiv.2512.22219}.}

\leavevmode\vadjust pre{\hypertarget{ref-ascend19716}{}}%
\CSLLeftMargin{{[}29{]} }%
\CSLRightInline{vLLM Ascend contributors. {Overlap Engram WKV
projections with MoE dispatch}.
\url{https://github.com/vllm-project/vllm-ascend/pull/19716}.}

\leavevmode\vadjust pre{\hypertarget{ref-ascend19909}{}}%
\CSLLeftMargin{{[}30{]} }%
\CSLRightInline{vLLM Ascend contributors. {Automatically budget FlashMLA
DCP overlap cores}.
\url{https://github.com/vllm-project/vllm-ascend/pull/19909}.}

\leavevmode\vadjust pre{\hypertarget{ref-kwon2020nimble}{}}%
\CSLLeftMargin{{[}31{]} }%
\CSLRightInline{W. Kwon, G.-I. Yu, E. Jeong, 和 B.-G. Chun. Nimble:
Lightweight and Parallel GPU Task Scheduling for Deep Learning. Advances
in Neural Information Processing Systems 33, 2020, 33, 8343--8354.
\url{https://proceedings.neurips.cc/paper/2020/hash/5f0ad4db43d8723d18169b2e4817a160-Abstract.html}.}

\leavevmode\vadjust pre{\hypertarget{ref-chen2025opara}{}}%
\CSLLeftMargin{{[}32{]} }%
\CSLRightInline{A. Chen, F. Xu, L. Han, 等. Opara: Exploiting Operator
Parallelism for Expediting DNN Inference on GPUs. IEEE Transactions on
Computers, 2025, 74, (1), 325--333. doi:
\href{https://doi.org/10.1109/TC.2024.3475589}{10.1109/TC.2024.3475589}.}

\leavevmode\vadjust pre{\hypertarget{ref-zhu2025nanoflow}{}}%
\CSLLeftMargin{{[}33{]} }%
\CSLRightInline{K. Zhu, Y. Gao, Y. Zhao, 等. {NanoFlow: Towards Optimal
Large Language Model Serving Throughput}. 19th USENIX Symposium on
Operating Systems Design and Implementation, 2025, 749--765.
\url{https://www.usenix.org/conference/osdi25/presentation/zhu-kan}.}

\leavevmode\vadjust pre{\hypertarget{ref-kamath2025pod}{}}%
\CSLLeftMargin{{[}34{]} }%
\CSLRightInline{A. K. Kamath, R. Prabhu, J. Mohan, 等. {POD-Attention:
Unlocking Full Prefill-Decode Overlap for Faster LLM Inference}.
Proceedings of the 30th ACM International Conference on Architectural
Support for Programming Languages and Operating Systems, Volume 2, 2025,
897--912. doi:
\href{https://doi.org/10.1145/3676641.3715996}{10.1145/3676641.3715996}.}

\leavevmode\vadjust pre{\hypertarget{ref-ma2020rammer}{}}%
\CSLLeftMargin{{[}35{]} }%
\CSLRightInline{L. Ma, Z. Xie, Z. Yang, 等. Rammer: Enabling Holistic
Deep Learning Compiler Optimizations with rTasks. 14th USENIX Symposium
on Operating Systems Design and Implementation (OSDI 20), 2020,
881--897.
\url{https://www.usenix.org/conference/osdi20/presentation/ma}.}

\leavevmode\vadjust pre{\hypertarget{ref-ascend21373}{}}%
\CSLLeftMargin{{[}36{]} }%
\CSLRightInline{vLLM Ascend contributors. {Prepare static kernels during
graph warmup}.
\url{https://github.com/vllm-project/vllm-ascend/pull/21373}.}

\leavevmode\vadjust pre{\hypertarget{ref-zeng2025medusa}{}}%
\CSLLeftMargin{{[}37{]} }%
\CSLRightInline{S. Zeng, M. Xie, S. Gao, 等. Medusa: Accelerating
Serverless LLM Inference with Materialization. Proceedings of the 30th
ACM International Conference on Architectural Support for Programming
Languages and Operating Systems, Volume 1, 2025, 653--668. doi:
\href{https://doi.org/10.1145/3669940.3707285}{10.1145/3669940.3707285}.}

\leavevmode\vadjust pre{\hypertarget{ref-yadav2026hybrid}{}}%
\CSLLeftMargin{{[}38{]} }%
\CSLRightInline{D. K. Yadav 和 T. Zhao. {Hybrid JIT-CUDA Graph
Optimization for Low-Latency Large Language Model Inference}. 2026 IEEE
50th Annual Computers, Software, and Applications Conference, 2026,
196--204. doi:
\href{https://doi.org/10.1109/COMPSAC69091.2026.00036}{10.1109/COMPSAC69091.2026.00036}.}

\leavevmode\vadjust pre{\hypertarget{ref-sglang2026refactor}{}}%
\CSLLeftMargin{{[}39{]} }%
\CSLRightInline{SGLang contributors. {SGLang refactor: Remove the NPU
graph runner}.
\url{https://github.com/sgl-project/sglang/commit/f96d5b1ca8a7817d0a514e81bdb80baf5bdaab8d}.}

\leavevmode\vadjust pre{\hypertarget{ref-nvidia2025cudaguide}{}}%
\CSLLeftMargin{{[}40{]} }%
\CSLRightInline{NVIDIA Corporation. CUDA C++ Programming Guide, Release
13.0 Update 1, CUDA Graphs. 2025.
\url{https://docs.nvidia.com/cuda/archive/13.0.1/cuda-c-programming-guide/index.html\#cuda-graphs}.}

\leavevmode\vadjust pre{\hypertarget{ref-cann910conditional}{}}%
\CSLLeftMargin{{[}41{]} }%
\CSLRightInline{Huawei. {aclmdlRIAddCondTask}.
\url{https://www.hiascend.com/doc_center/source/en/CANNCommunityEdition/910/API/runtimeapi/aclcppdevg_03_2270.html}.}

\leavevmode\vadjust pre{\hypertarget{ref-audit_cann_cond_params}{}}%
\CSLLeftMargin{{[}42{]} }%
\CSLRightInline{Huawei. {aclmdlRICondTaskParams}.
\url{https://www.hiascend.com/doc_center/source/en/CANNCommunityEdition/910/API/runtimeapi/aclcppdevg_03_2274.html}.}

\leavevmode\vadjust pre{\hypertarget{ref-audit_cann_cond_create}{}}%
\CSLLeftMargin{{[}43{]} }%
\CSLRightInline{Huawei. {aclmdlRICondHandleCreate}.
\url{https://www.hiascend.com/doc_center/source/en/CANNCommunityEdition/910/API/runtimeapi/aclcppdevg_03_2268.html}.}

\leavevmode\vadjust pre{\hypertarget{ref-sglang8030}{}}%
\CSLLeftMargin{{[}44{]} }%
\CSLRightInline{SGLang contributors. {SGLang issue8030: support
ACLGraph}. \url{https://github.com/sgl-project/sglang/issues/8030}.}

\leavevmode\vadjust pre{\hypertarget{ref-sglang12194}{}}%
\CSLLeftMargin{{[}45{]} }%
\CSLRightInline{SGLang contributors. {SGLang issue12194: ACLGraph
Instructions And Precaution}.
\url{https://github.com/sgl-project/sglang/issues/12194}.}

\leavevmode\vadjust pre{\hypertarget{ref-vllmcudagraphdesign}{}}%
\CSLLeftMargin{{[}46{]} }%
\CSLRightInline{vLLM contributors. {CUDA Graphs design}.
\url{https://docs.vllm.ai/en/v0.11.0/design/cuda_graphs.html}.}

\end{CSLReferences}

\end{document}
